% Ampliación estadística para el artículo V (ES).
% Fuente principal: N19; antecedente de memoria ordenada: F044.
% Este archivo no modifica la fuente activa y debe integrarse como sucesor.
\section{Dualité de Fourier, incertitude finie et mémoire ordonnée}
\label{v:n19:sec:incertidumbre-fourier}

Le domaine de la réalisation finie provient de la chaîne typée déjà
construite dans l’article :
\begin{equation}
 \mathsf{APP}\longrightarrow\mathsf{TRIT}\longrightarrow\mathsf{TPK}
 \longrightarrow X_{\mathrm{enr}}
 \longrightarrow (\mathbb Z/9\mathbb Z)^3
 \xrightarrow{\,\beta_{729}\,}
 \Lambda:=(\mathbb Z/27\mathbb Z)^2
 \longrightarrow \ell^2(\Lambda).
 \label{v:n19:eq:cadena-tipificada}
\end{equation}
Ici \(X_{\mathrm{enr}}\) conserve le chemin, l’orientation, la retenue, le bord et la
mémoire ; le passage à \((\mathbb Z/9\mathbb Z)^3\) est le registre positionnel
pertinent. L’application \(\beta_{729}\), définie dans
\eqref{v:ant:eq:ii-area-beta729} et démontrée bijective dans la
Proposition~\ref{v:ant:prop:ii-area-beta729}, regroupe les six trits
ordonnés en deux coordonnées d’ordre vingt-sept. Le registre positionnel
et la réalisation de groupe sont ainsi distingués : \(\beta_{729}\)
transporte l’incidence ordonnée, tandis que la transformée de
Fourier emploie les opérations propres de
\((\mathbb Z/27\mathbb Z)^2\). La généalogie détermine donc quelque chose de
plus précis que le cardinal commun de \(729\) positions.

\subsection{Caractères unitaires et bases mutuellement non biaisées}
\label{v:n19:sec:caracteres-mub}

Écrivons \(N:=|\Lambda|=27^2=729\). Pour
\(x=(x_1,x_2)\) et \(k=(k_1,k_2)\) dans \(\Lambda\), soit
\begin{equation}
 \chi_k(x):=
 \exp\!\left(\frac{2\pi i}{27}(k_1x_1+k_2x_2)\right).
 \label{v:n19:eq:caracteres}
\end{equation}
Le produit scalaire des coordonnées est pris modulo vingt-sept.
L’orthogonalité des caractères donne
\begin{equation}
 \sum_{x\in\Lambda}\chi_k(x)\overline{\chi_{\ell}(x)}
 =N\,\delta_{k,\ell}.
 \label{v:n19:eq:ortogonalidad-caracteres}
\end{equation}
Nous définissons la transformée unitaire \(\mathcal F_\Lambda\) par
\begin{equation}
 \widehat\psi(k):=(\mathcal F_\Lambda\psi)(k)
 =\frac1{\sqrt N}\sum_{x\in\Lambda}\chi_k(x)\psi(x),
 \qquad
 \psi(x)=\frac1{\sqrt N}\sum_{k\in\Lambda}
 \overline{\chi_k(x)}\widehat\psi(k).
 \label{v:n19:eq:fourier-inversion}
\end{equation}
Le caractère conjugué de la seconde formule est indispensable sous la
convention positive de la première. L’identité
\eqref{v:n19:eq:ortogonalidad-caracteres} prouve simultanément
l’inversion et \(\|\mathcal F_\Lambda\psi\|_2=\|\psi\|_2\).

\begin{lemma}[Absence de biais mutuel des bases]
\label{v:n19:lem:mub}
Soient \(\{e_x:x\in\Lambda\}\) la base de position et
\(f_k(x)=N^{-1/2}\overline{\chi_k(x)}\) la base de Fourier. Alors
\begin{equation}
 |\langle e_x,f_k\rangle|=N^{-1/2}=1/27
 \qquad(x,k\in\Lambda).
 \label{v:n19:eq:mub}
\end{equation}
Par conséquent, les deux bases sont mutuellement non biaisées.
\end{lemma}
\begin{proof}
L’orthonormalité de \(\{f_k\}\) découle de
\eqref{v:n19:eq:ortogonalidad-caracteres} ; en outre,
\(\langle e_x,f_k\rangle=f_k(x)\) et tout caractère est de module un.
\end{proof}

\subsection{Incertitude exacte de support}
\label{v:n19:sec:soporte}

Pour une fonction \(u:\Lambda\to\mathbb C\), notons
\(\operatorname{supp}u=\{x:u(x)\ne0\}\) son support.

\begin{theorem}[Borne de support dans le tore triadique]
\label{v:n19:thm:soporte}
Pour tout \(\psi\in\ell^2(\Lambda)\) non nul,
\begin{equation}
 |\operatorname{supp}\psi|\,
 |\operatorname{supp}\widehat\psi|\ge N=729.
 \label{v:n19:eq:cota-soporte}
\end{equation}
\end{theorem}
\begin{proof}
Posons \(S=\operatorname{supp}\psi\) et
\(T=\operatorname{supp}\widehat\psi\). L’inversion correcte de
\eqref{v:n19:eq:fourier-inversion} donne, pour tout \(x\in S\),
\[
 |\psi(x)|
 =\left|\frac1{\sqrt N}\sum_{k\in T}
       \overline{\chi_k(x)}\widehat\psi(k)\right|
 \le \sqrt{\frac{|T|}{N}}\,\|\widehat\psi\|_2,
\]
par Cauchy--Schwarz et \(|\chi_k(x)|=1\). En élevant au carré et
en sommant sur \(S\),
\[
 \|\psi\|_2^2
 \le\frac{|S||T|}{N}\,\|\widehat\psi\|_2^2
 =\frac{|S||T|}{N}\,\|\psi\|_2^2.
\]
La norme est positive ; sa simplification prouve
\eqref{v:n19:eq:cota-soporte}. C’est la spécialisation au groupe
fini réalisé de l’incertitude de support de
Donoho--Stark~\cite{donoho_stark_1989}.
\end{proof}

La chaîne \eqref{v:n19:eq:cadena-tipificada} fixe le domaine avant de
formuler l’inégalité : APP--TRIT--TPK apporte le registre positionnel et
\(\beta_{729}\) sa réalisation finie comme groupe \(\Lambda\).

\subsection{Annulateurs et classification des cas de saturation}
\label{v:n19:sec:anuladores}

Si \(H\le\Lambda\), son annulateur est
\begin{equation}
 H^\perp:=\{k\in\Lambda:\chi_k(h)=1\text{ pour tout }h\in H\}.
 \label{v:n19:eq:anulador}
\end{equation}
Pour des diviseurs \(a,b\mid27\), considérons
\begin{equation}
 H_{a,b}:=\{(x_1,x_2)\in\Lambda:
 x_1\equiv0\pmod a,\ x_2\equiv0\pmod b\}.
 \label{v:n19:eq:subgrupo-rectangular}
\end{equation}
Alors
\begin{equation}
 |H_{a,b}|=\frac{27}{a}\frac{27}{b},
 \quad
 H_{a,b}^{\perp}
 =\{(k_1,k_2):k_1\equiv0\pmod{27/a},
                  k_2\equiv0\pmod{27/b}\},
 \quad |H_{a,b}^{\perp}|=ab.
 \label{v:n19:eq:anulador-explicito}
\end{equation}
En effet, la somme géométrique dans chaque coordonnée est nulle sauf lorsque
\(27\mid ak_1\) et \(27\mid bk_2\). Par conséquent, pour l’indicatrice
normalisée
\begin{equation}
 \psi_{a,b}:=|H_{a,b}|^{-1/2}\mathbf1_{H_{a,b}},
 \label{v:n19:eq:indicador-normalizado}
\end{equation}
on obtient exactement
\begin{equation}
 \widehat\psi_{a,b}(k)
 =\sqrt{\frac{|H_{a,b}|}{N}}\,
   \mathbf1_{H_{a,b}^{\perp}}(k).
 \label{v:n19:eq:fourier-indicador}
\end{equation}
Ainsi, la transformée est uniforme sur l’annulateur et
\(|H_{a,b}|\,|H_{a,b}^{\perp}|=N\).

\begin{table}[htbp]
\centering
\caption{Quatre saturations exactes, avec domaines différenciés.}
\label{v:n19:tab:saturaciones}
\begin{tabular}{@{}ccccc@{}}
\toprule
\((a,b)\) & \(|H_{a,b}|\) & \(|H_{a,b}^{\perp}|\) & produit
 & \((H_{\rm pos},H_{\rm F})\) \\
\midrule
\((3,3)\)  & \(81\) & \(9\)  & \(729\) & \((\log81,\log9)\) \\
\((9,9)\)  & \(9\)  & \(81\) & \(729\) & \((\log9,\log81)\) \\
\((3,9)\)  & \(27\) & \(27\) & \(729\) & \((\log27,\log27)\) \\
\((1,27)\) & \(27\) & \(27\) & \(729\) & \((\log27,\log27)\) \\
\bottomrule
\end{tabular}
\end{table}
Les deux dernières lignes ont les mêmes cardinaux, mais ni le
même sous-groupe ni le même annulateur ; l’égalité numérique ne confond pas leurs
incidences.

\subsection{Incertitude entropique par interpolation et différentiation}
\label{v:n19:sec:entropia}

Soit \(\|\psi\|_2=1\), et définissons les distributions
\begin{equation}
 p_x:=|\psi(x)|^2,\qquad q_k:=|\widehat\psi(k)|^2,
 \qquad
 H_{\rm pos}:=-\sum_xp_x\log p_x,\quad
 H_{\rm F}:=-\sum_kq_k\log q_k,
 \label{v:n19:eq:entropias-distribucion}
\end{equation}
avec logarithme naturel et convention \(0\log0=0\). Ce sont les entropies de
Shannon de deux distributions de mesure sur des bases conjuguées.

\begin{theorem}[Borne entropique de la réalisation \(\Lambda\)]
\label{v:n19:thm:entropia}
Tout état unitaire satisfait
\begin{equation}
 H_{\rm pos}(\psi)+H_{\rm F}(\psi)
 \ge\log N=\log729=6\log3.
 \label{v:n19:eq:cota-entropica}
\end{equation}
\end{theorem}
\begin{proof}
De \(\|\mathcal F_\Lambda\|_{\ell^1\to\ell^\infty}
\le N^{-1/2}\) et de l’unitarité
\(\|\mathcal F_\Lambda\|_{\ell^2\to\ell^2}=1\), le théorème
d’interpolation de Riesz--Thorin~\cite{riesz_1927,thorin_1948} donne, pour
\(1<p\le2\) et \(p'=p/(p-1)\),
\begin{equation}
 \|\mathcal F_\Lambda\psi\|_{p'}
 \le N^{\frac12-\frac1p}\|\psi\|_p.
 \label{v:n19:eq:riesz-thorin}
\end{equation}
Considérons
\begin{equation}
 g(p):=\log\|\mathcal F_\Lambda\psi\|_{p'}
       -\log\|\psi\|_p
       -\left(\frac12-\frac1p\right)\log N.
 \label{v:n19:eq:funcion-interpolacion}
\end{equation}
L’inégalité \eqref{v:n19:eq:riesz-thorin} implique \(g(p)\le0\),
tandis que \(g(2)=0\) ; c’est pourquoi \(g'_-(2)\ge0\).

Pour tout vecteur unitaire \(u\) en dimension finie,
\begin{equation}
 \left.\frac{d}{dr}\log\|u\|_r\right|_{r=2}
 =\frac14\sum_j|u_j|^2\log|u_j|^2
 =-\frac14H(|u|^2).
 \label{v:n19:eq:derivada-norma}
\end{equation}
Les termes nuls s’obtiennent par continuité ; les termes positifs
sont différentiables. Comme \(\left.dp'/dp\right|_{p=2}=-1\), en dérivant
\eqref{v:n19:eq:funcion-interpolacion} à gauche, on obtient
\begin{equation}
 g'_-(2)=\frac14
 \bigl(H_{\rm pos}+H_{\rm F}-\log N\bigr)\ge0,
 \label{v:n19:eq:derivada-entropica}
\end{equation}
ce qui équivaut à \eqref{v:n19:eq:cota-entropica}.

Cette dérivation s’inscrit dans la lignée de l’incertitude entropique de
Hirschman et de Beckner~\cite{hirschman_1957,beckner_1975} ; pour deux bases
arbitraires, la formulation par recouvrement maximal est celle de
Maassen--Uffink~\cite{maassen_uffink_1988}. Dans la présente réalisation,
la constante \(\log N\) s’obtient directement en différentiant le facteur
d’interpolation de \eqref{v:n19:eq:riesz-thorin}.
\end{proof}

Pour les états \(\psi_{a,b}\), \eqref{v:n19:eq:fourier-indicador}
montre que \(p\) et \(q\) sont uniformes sur \(H_{a,b}\) et
\(H_{a,b}^{\perp}\), respectivement. Par conséquent,
\[
 H_{\rm pos}=\log|H_{a,b}|,\qquad
 H_{\rm F}=\log|H_{a,b}^{\perp}|,\qquad
 H_{\rm pos}+H_{\rm F}=\log729,
\]
ce qui prouve la saturation entropique des quatre cas du
Tableau~\ref{v:n19:tab:saturaciones} sans approximation numérique.

\subsection{Domaines typés de l’information}
\label{v:n19:sec:tipos-entropia}

La quantité \(H_{\rm pos}+H_{\rm F}\) mesure la somme de deux entropies
classiques de distribution associées à des mesures conjuguées. L’entropie
d’intrication de la coupure de bord, l’entropie de von Neumann d’un
état réduit, l’entropie thermodynamique et la température, et l’entropie
marginale de l’histogramme tritique sont définies sur des objets et des domaines
distincts. Toute comparaison entre ces grandeurs est formulée au moyen d’une
application explicite qui spécifie le domaine, le codomaine et l’état. La borne
\eqref{v:n19:eq:cota-entropica} appartient intégralement à
\(\ell^2(\Lambda)\).

\subsection{Mémoire ordonnée dans le domaine dodécaphasique}
\label{v:n19:sec:memoria-ordenada}

Le résultat F044 agit sur des mots dodécaphasiques et caractérise
l’information d’ordre conservée par l’état enrichi. Son domaine et
ses opérateurs sont indépendants de la réalisation de Fourier de N19 ;
la relation entre les deux résultats compare l’information marginale à
l’information ordonnée.

Soit \(\tau=(\tau_1,\ldots,\tau_{12})\in\{-1,0,+1\}^{12}\), avec
\(\operatorname{sgn}(0)=0\), et des indices cycliques modulo douze. Définissons
\begin{align}
 G_a&:=\{a,a+4,a+8\},\qquad a=1,2,3,4,\nonumber\\
 \nu_{120}(\tau)&:=\sum_{a=1}^4
   \operatorname{sgn}\!\left(\sum_{m\in G_a}\tau_m\right),
 \label{v:n19:eq:nu120}\\
 \nu_{270}(\tau)&:=\sum_{m=1}^{12}\tau_m\tau_{m+3}.
 \label{v:n19:eq:nu270}
\end{align}
Les mots
\begin{equation}
 \tau^{(a)}=(1,1,0,0,0,0,0,0,0,0,0,0),\qquad
 \tau^{(b)}=(1,0,0,1,0,0,0,0,0,0,0,0)
 \label{v:n19:eq:testigos-orden}
\end{equation}
ont le même histogramme ---dix zéros et deux unités--- et donc
la même entropie marginale d’alphabet. Toutefois,
\begin{equation}
 (\nu_{120},\nu_{270})(\tau^{(a)})=(2,0),\qquad
 (\nu_{120},\nu_{270})(\tau^{(b)})=(2,1).
 \label{v:n19:eq:testigos-lectura}
\end{equation}
Le second lecteur détecte une incidence à distance trois que
l’histogramme élimine.

Pour l’involution \(C_M\tau:=-\operatorname{rev}(\tau)\), le renversement
permute les quatre classes \(G_a\) et la négation inverse leurs sommes ;
l’autocorrélation cyclique de pas \(+3\) coïncide, après réindexation, avec
celle de pas \(-3\). Par conséquent,
\begin{equation}
 \nu_{120}(C_M\tau)=-\nu_{120}(\tau),\qquad
 \nu_{270}(C_M\tau)=\nu_{270}(\tau).
 \label{v:n19:eq:paridades-lectores}
\end{equation}
Le lecteur \(\nu_{120}\) est impair et le lecteur quadratique \(\nu_{270}\) est
pair. Ces parités caractérisent l’information supplémentaire que conserve
le mot ordonné.

\subsection{Domaine de validité et conditions mathématiques}
\label{v:n19:sec:falsadores}

Les conditions qui délimitent les deux résultats sont :
\begin{enumerate}
 \item \(\beta_{729}\) est bijective sur la réalisation positionnelle de
 \(\Lambda\) utilisée dans \eqref{v:n19:eq:cadena-tipificada} ;
 \item l’orthogonalité \eqref{v:n19:eq:ortogonalidad-caracteres} et
 l’inversion avec caractère conjugué établissent l’unitarité ;
 \item tout état non nul satisfait le produit des supports
 \(\lvert\operatorname{supp}\psi\rvert
 \lvert\operatorname{supp}\widehat\psi\rvert\ge729\) ;
 \item tout état unitaire satisfait
 \(H_{\rm pos}+H_{\rm F}\ge\log729\) ;
 \item dans les quatre cas du
 Tableau~\ref{v:n19:tab:saturaciones}, la transformée est uniforme sur
 l’annulateur et le produit des cardinaux est \(729\) ;
 \item les témoins de \eqref{v:n19:eq:testigos-orden} partagent
 l’histogramme, produisent \eqref{v:n19:eq:testigos-lectura} et satisfont les
 parités \eqref{v:n19:eq:paridades-lectores}.
\end{enumerate}
Les cinq premiers énoncés appartiennent à la réalisation finie typée
\(\Lambda\) et à sa transformée de Fourier unitaire ; le sixième appartient au domaine
dodécaphasique de mémoire ordonnée. Ces conditions fournissent aussi
des critères de réfutation locaux : le non-respect d’une identité, d’une inégalité ou d’une
valeur témoin réfute l’énoncé correspondant dans son domaine. La
relation entre N19 et F044 est complémentaire et maintient séparés leurs
domaines et opérateurs.
